\documentclass[a4paper,12pt]{article}
\usepackage[utf8]{inputenc}
\usepackage[russian]{babel}
\usepackage{amsmath,amssymb}
\usepackage{geometry}
\usepackage{graphicx} % Пакет для картинок

\geometry{margin=2cm}

\begin{document}

\section*{КР3 }

\subsection*{1) МП-автомат для $(abc)^n(cba)^k(ad)^m(dc)^p:\; n+m=k+p,\; n\ge k$}

\textbf{Краткое обоснование:} стек хранит разность $(n-k)+(m-p)$: на каждом блоке $abc$ и $ad$ делаем push(1), на каждом $cba$ и $dc$ делаем pop(1); условие $n\ge k$ гарантирует, что на этапе $cba$ не потребуется снимать с пустого стека. Принимаем, если стек пуст.

\begin{figure}[h]
    \centering
    \includegraphics[width=0.9\textwidth]{photo_2026-01-13_16-18-27.jpg} 
\caption { автоматы делать в латехе очень сложно я не смог посему как есть}
\end{figure}

\newpage
\subsection*{2) ДКА для регулярного выражения $a?ba?b^*ab$}

\textbf{Краткое обоснование:} ДКА построен по стандартному алгоритму (через НКА и детерминизацию). Он принимает строки, соответствующие шаблону: необязательное $a$, затем $b$, затем необязательное $a$, любое кол-во $b$, и в конце $ab$.

\begin{figure}[h]
    \centering
    \includegraphics[width=0.9\textwidth]{photo_2026-01-13_16-18-40.jpg}
    \caption{аналогично, смог только вручную}
\end{figure}

\newpage
\subsection*{3) Регулярное выражение по данному ДКА}


\begin{figure}[h]
    \centering
    \includegraphics[width=0.9\textwidth]{photo_2026-01-13_16-29-35.jpg}
    \caption{регулярки тоже непросто в латехе чтобы читаемы были}
\end{figure}

\end{document}
